% \chapter{Introduction, Background, and Machine Learning Problem}
% \label{ch:introduction}

% This formative report documents the CMP6230 data management and analytics work for the Lichess Open Database: planning and design (Task~1), followed by implementation of an MLOps pipeline (Task~2). The module learning outcome is to model and communicate data management requirements for analytical operations---storage, retrieval, validation, and downstream machine learning---with justified structure and tooling at each stage.

% Chess platforms generate large volumes of rated game telemetry. The supervised learning problem addressed here is to predict game outcomes (or related rating-based signals) from pre-game headers, openings, and player metadata embedded in Portable Game Notation (PGN) archives. That task is non-trivial at Lichess scale: monthly shards are multi-gigabyte compressed blobs, semi-structured rather than relational, and refreshed on a public schedule. A credible pipeline must therefore support both ad hoc SQL and notebook analytics and an automated training path without corrupting temporal structure.

% Chapter~\ref{ch:candidate-datasets} compares three candidate datasets and justifies Lichess as the corpus for this project. Chapter~\ref{ch:storage-planning} details the initial data storage plan. Chapter~\ref{ch:initial-pipeline-planning} presents the \emph{initial} end-to-end MLOps design from Task~1. Chapter~\ref{ch:final-pipeline-planning} refines that design into the \emph{final} pipeline plan informed by implementation. Subsequent chapters describe the runnable system, reflective discussion, and plans for the summative submission.
